\textbf{Background/Objectives:} Small multiple sclerosis (MS) lesions are easily missed on MRI, yet the 2024 McDonald criteria consider lesions from 3~mm. Models trained for voxel overlap (Dice) can score well while missing them. \textbf{Methods:} CATMIL adds two terms to the nnU-Net loss: a multiple instance learning (MIL) term penalizing completely missed lesions and a component-adaptive Tversky (CAT) term weighting every lesion equally. Models using T1-weighted, T2-weighted, and FLAIR MRI from two public datasets (MSLesSeg, 53 patients; 3D-MR-MS, 30 patients) were compared with four alternative losses. \textbf{Results:} For lesions $\geq$3~mm, recall was 94.4\% versus 87.0\% with the standard loss on MSLesSeg (95\% CI +3.2 to +12.9 points) and 89.4\% versus 83.0\% on 3D-MR-MS. On MSLesSeg, missed lesions fell from 5.0 to 2.6 per scan with similar Dice, reproduced in 47 cross-validation patients; per model, CATMIL found 30.6 lesions the standard loss missed, versus 2.0 the reverse. In an ablation, MIL drove detection; adding CAT gave the best Dice and boundary accuracy. The cost was more small false positives, mostly removed by a 5-voxel filter on MSLesSeg but larger on 3D-MR-MS (lesion-wise F1 0.47 vs.\ 0.60--0.67; 20-voxel filter needed). \textbf{Conclusions:} CATMIL substantially reduces missed small MS lesions, supporting its use as a second reader, at a dataset-dependent false-positive cost.
